% GENERATED FILE. Edit canonical lessons, then run npm run generate:books.
% canonical-source-hash: fnv1a64:94fa9c6db7d841b7
% canonical-lessons: JA-C528-aruiwa, JA-C528-matawa, JA-C528-soretomo, JA-C528-moshikuwa, JA-C528-izureka

\chapter{This or That}
\label{ch:ja-this-or-that}
\hlchaptermodality{528}

\begin{chapteropening}
\textbf{By the end of this chapter:} \emph{I can offer a choice between two things.}

Complete the last lesson of chapter 528: I can offer a choice between two things.
\end{chapteropening}

\section[\texorpdfstring{\ja{あるいは} (aruiwa)}{あるいは (aruiwa)}]{\ja{あるいは} (aruiwa) — or perhaps; or}

\label{lesson:JA-C528-aruiwa}



\begin{quote}
Before the new one: say the Japanese for from now on, then the Japanese for the present day.
\end{quote}

\subsection*{You'll want to know: \ja{あるいは}}
\textbf{\ja{あるいは}} — \emph{aruiwa} — \textquotedblleft{}or perhaps; or\textquotedblright{}.

\begin{quote}
\textbf{\ja{あした} \ja{さむく} \ja{なったら、あめ} \ja{あるいは} \ja{ゆきでしょう。}} — \emph{ashita samuku nattara, ame aruiwa yuki deshō} — \textquotedblleft{}If it gets cold tomorrow, there will be rain or snow.\textquotedblright{}
\end{quote}

\begin{cousinweb}
\textbf{\ja{あるいは}} (\emph{aruiwa}) is a native Japanese word (wago), inherited rather than borrowed from Chinese, and usually written in kana; its older origin is not settled (JMdict).
\end{cousinweb}

\begin{grammarlens}[title={What you've built}]
One more word for saying what happens if.
\end{grammarlens}

\subsection*{Your turn}
\begin{itemize}
\raggedright
  \item \emph{Say it:} \emph{aruiwa}
  \item \emph{Say it:} \emph{aruiwa}, once more
  \item \emph{Say it:} \emph{ashita samuku nattara, ame aruiwa yuki deshō} — If it gets cold tomorrow, there will be rain or snow.
  \item \emph{Recall:} say the Japanese for from now on, then the Japanese for the present day, then say \emph{aruiwa} again
\end{itemize}

\begin{tcolorbox}[breakable,colback=teal!4,colframe=teal!35!black,title={Before you move on}]
What does \ja{あるいは} mean? (\textquotedblleft{}Or perhaps; or\textquotedblright{}.) Say it once more.
\end{tcolorbox}
\section[\texorpdfstring{\ja{または} (matawa)}{または (matawa)}]{\ja{または} (matawa) — or; or else}

\label{lesson:JA-C528-matawa}



\begin{quote}
Before the new one: say the Japanese for the present day, then the Japanese for or perhaps.
\end{quote}

\subsection*{You'll want to know: \ja{または}}
\textbf{\ja{または}} — \emph{matawa} — \textquotedblleft{}or; or else\textquotedblright{}.

\begin{quote}
\textbf{\ja{もし} \ja{あめが} \ja{ふったら、でんしゃ} \ja{または} \ja{くるまで} \ja{いきます。}} — \emph{moshi ame ga futtara, densha matawa kuruma de ikimasu} — \textquotedblleft{}If it rains, I will go by train or by car.\textquotedblright{}
\end{quote}

\begin{cousinweb}
\textbf{\ja{または}} (\emph{matawa}) is a native Japanese word (wago), inherited rather than borrowed from Chinese, and usually written in kana; its older origin is not settled (JMdict).
\end{cousinweb}

\begin{grammarlens}[title={What you've built}]
One more word for saying what happens if.
\end{grammarlens}

\subsection*{Your turn}
\begin{itemize}
\raggedright
  \item \emph{Say it:} \emph{matawa}
  \item \emph{Say it:} \emph{matawa}, once more
  \item \emph{Say it:} \emph{moshi ame ga futtara, densha matawa kuruma de ikimasu} — If it rains, I will go by train or by car.
  \item \emph{Recall:} say the Japanese for the present day, then the Japanese for or perhaps, then say \emph{matawa} again
\end{itemize}

\begin{tcolorbox}[breakable,colback=teal!4,colframe=teal!35!black,title={Before you move on}]
What does \ja{または} mean? (\textquotedblleft{}Or; or else\textquotedblright{}.) Say it once more.
\end{tcolorbox}
\section[\texorpdfstring{\ja{それとも} (soretomo)}{それとも (soretomo)}]{\ja{それとも} (soretomo) — or (in a question)}

\label{lesson:JA-C528-soretomo}



\begin{quote}
Before the new one: say the Japanese for or perhaps, then the Japanese for or.
\end{quote}

\subsection*{You'll want to know: \ja{それとも}}
\textbf{\ja{それとも}} — \emph{soretomo} — \textquotedblleft{}or (in a question)\textquotedblright{}.

\begin{quote}
\textbf{\ja{あめが} \ja{ふったら、えいがに} \ja{しますか、それとも} \ja{うちに} \ja{いますか。}} — \emph{ame ga futtara, eiga ni shimasu ka, soretomo uchi ni imasu ka} — \textquotedblleft{}If it rains, shall we see a film, or stay at home?\textquotedblright{}
\end{quote}

\begin{cousinweb}
\textbf{\ja{それとも}} (\emph{soretomo}) is a native Japanese word (wago), inherited rather than borrowed from Chinese, and usually written in kana; its older origin is not settled (JMdict).
\end{cousinweb}

\begin{grammarlens}[title={What you've built}]
One more word for saying what happens if.
\end{grammarlens}

\subsection*{Your turn}
\begin{itemize}
\raggedright
  \item \emph{Say it:} \emph{soretomo}
  \item \emph{Say it:} \emph{soretomo}, once more
  \item \emph{Say it:} \emph{ame ga futtara, eiga ni shimasu ka, soretomo uchi ni imasu ka} — If it rains, shall we see a film, or stay at home?
  \item \emph{Recall:} say the Japanese for or perhaps, then the Japanese for or, then say \emph{soretomo} again
\end{itemize}

\begin{tcolorbox}[breakable,colback=teal!4,colframe=teal!35!black,title={Before you move on}]
What does \ja{それとも} mean? (\textquotedblleft{}Or (in a question)\textquotedblright{}.) Say it once more.
\end{tcolorbox}
\section[\texorpdfstring{\ja{もしくは} (moshikuwa)}{もしくは (moshikuwa)}]{\ja{もしくは} (moshikuwa) — or else; or}

\label{lesson:JA-C528-moshikuwa}



\begin{quote}
Before the new one: say the Japanese for or, then the Japanese for or (in a question).
\end{quote}

\subsection*{You'll want to know: \ja{もしくは}}
\textbf{\ja{もしくは}} — \emph{moshikuwa} — \textquotedblleft{}or else; or\textquotedblright{}.

\begin{quote}
\textbf{\ja{もし} \ja{あめが} \ja{ふったら、しあいは} \ja{げつようび} \ja{もしくは} \ja{かようびです。}} — \emph{moshi ame ga futtara, shiai wa getsuyōbi moshikuwa kayōbi desu} — \textquotedblleft{}If it rains, the match will be on Monday or Tuesday.\textquotedblright{}
\end{quote}

\begin{cousinweb}
\textbf{\ja{もしくは}} (\emph{moshikuwa}) is a native Japanese word (wago), inherited rather than borrowed from Chinese, and usually written in kana; its older origin is not settled (JMdict).
\end{cousinweb}

\begin{grammarlens}[title={What you've built}]
One more word for saying what happens if.
\end{grammarlens}

\subsection*{Your turn}
\begin{itemize}
\raggedright
  \item \emph{Say it:} \emph{moshikuwa}
  \item \emph{Say it:} \emph{moshikuwa}, once more
  \item \emph{Say it:} \emph{moshi ame ga futtara, shiai wa getsuyōbi moshikuwa kayōbi desu} — If it rains, the match will be on Monday or Tuesday.
  \item \emph{Recall:} say the Japanese for or, then the Japanese for or (in a question), then say \emph{moshikuwa} again
\end{itemize}

\begin{tcolorbox}[breakable,colback=teal!4,colframe=teal!35!black,title={Before you move on}]
What does \ja{もしくは} mean? (\textquotedblleft{}Or else; or\textquotedblright{}.) Say it once more.
\end{tcolorbox}
\section[\texorpdfstring{\ja{いずれか} (izureka)}{いずれか (izureka)}]{\ja{いずれか} (izureka) — either one; one of them}

\label{lesson:JA-C528-izureka}



\begin{quote}
Before the new one: say the Japanese for or (in a question), then the Japanese for or else.
\end{quote}

\subsection*{You'll want to know: \ja{いずれか}}
\textbf{\ja{いずれか}} — \emph{izureka} — \textquotedblleft{}either one; one of them\textquotedblright{}.

\begin{quote}
\textbf{\ja{もし} \ja{じかんが} \ja{あったら、いずれかの} \ja{しあいに} \ja{きて} \ja{ください。}} — \emph{moshi jikan ga attara, izureka no shiai ni kite kudasai} — \textquotedblleft{}If you have time, please come to one of the matches.\textquotedblright{}
\end{quote}

\begin{cousinweb}
\textbf{\ja{いずれか}} (\emph{izureka}) is a native Japanese word (wago), inherited rather than borrowed from Chinese, and usually written in kana: built from \emph{izure}, \textquotedblleft{}which\textquotedblright{}, and the particle \emph{ka} (JMdict).
\end{cousinweb}

\begin{grammarlens}[title={What you've built}]
One more word for saying what happens if.
\end{grammarlens}

\subsection*{Your turn}
\begin{itemize}
\raggedright
  \item \emph{Say it:} \emph{izureka}
  \item \emph{Say it:} \emph{izureka}, once more
  \item \emph{Say it:} \emph{moshi jikan ga attara, izureka no shiai ni kite kudasai} — If you have time, please come to one of the matches.
  \item \emph{Recall:} say the Japanese for or (in a question), then the Japanese for or else, then say \emph{izureka} again
\end{itemize}

\begin{tcolorbox}[breakable,colback=teal!4,colframe=teal!35!black,title={Before you move on}]
What does \ja{いずれか} mean? (\textquotedblleft{}Either one; one of them\textquotedblright{}.) Say it once more.
\end{tcolorbox}
